% ---
% Subsección 2.1.9 de los Antecedentes: versión ampliada de la tabla tab:estado_arte
% (líneas 831–974 de Docs/Technical_Report.tex). Va al final del capítulo porque resume los trabajos
% de 2.1.1–2.1.8 y usa las cuatro preguntas de la sección 2.1.8.
% Cambios respecto a TT1: sin las filas de Bilal, Singh y Ajlouni (sin fuente verificable; decisión del
% autor del 2026-10-16); columnas P1–P4 en lugar de una valoración agregada; resultados del presente
% trabajo (H-042) en lugar de «Por det.».
% ---

\subsubsection{Tabla comparativa y brechas identificadas}

La tabla~\ref{tab:estado_arte} resume los trabajos revisados en este capítulo. Las columnas P1 a P4 recogen la respuesta de cada trabajo a las cuatro preguntas de control de la sección~\ref{sec:revision_critica} (paridad, incertidumbre, partición de prueba y simulabilidad); no se agregan en una sola valoración, porque responden a cuestiones distintas y la cuarta no la satisface ningún experimento a esta escala.

Entre los trabajos revisados, ninguno reúne a la vez un circuito cuántico de codificación, el CBIS-DDSM como conjunto de datos, un análisis explícito de la separabilidad de clases en el espacio inducido por el circuito y el tratamiento de masas y calcificaciones como subproblemas separados. Esa combinación define el lugar del presente trabajo. Las dos últimas condiciones corresponden a las líneas de trabajo futuro que señalan Azevedo et al. (sección~\ref{sec:trabajos_cbis_ddsm}).

\begin{landscape}
\begin{table}[htbp]
\centering
\caption{Comparativa de trabajos relacionados en el estado del arte}
\label{tab:estado_arte_v2}
\label{tab:estado_arte}
\renewcommand{\arraystretch}{1.08}
\setlength{\tabcolsep}{3pt}
\scriptsize

\newcolumntype{L}[1]{>{\raggedright\arraybackslash}p{#1}}
\begin{tabularx}{\linewidth}{@{}L{2.4cm}L{2.5cm}L{1.4cm}L{2.3cm}L{2.3cm}L{1.1cm}L{1.2cm}L{1.3cm}*{4}{L{0.7cm}}>{\raggedright\arraybackslash}X@{}}

\toprule
\textbf{Trabajo} &
\textbf{Datos} &
\textbf{Circuito cuántico} &
\textbf{Modelo / framework} &
\textbf{Feature map} &
\textbf{CBIS-DDSM} &
\textbf{Análisis separab.} &
\textbf{Masas/ calcs. sep.} &
\textbf{P1} & \textbf{P2} & \textbf{P3} & \textbf{P4} &
\textbf{Desempeño reportado} \\
\midrule

\multicolumn{13}{@{}l}{\cellcolor{gray!12}\scriptsize\textit{Aprendizaje profundo en mamografía}} \\[1pt]

Shen et al. (2019) \cite{Shen2019DeepLearning}
& CBIS-DDSM
& \xmark
& ResNet / VGG
& ---
& \cmark
& \xmark
& \xmark
& --- & --- & --- & ---
& AUC: 0.91 \\

Wang et al. (2020) \cite{Wang2020DeepLearning}
& DDSM / INbreast / MIAS / UKy
& \xmark
& CNN (6 modelos)
& ---
& Parcial
& \xmark
& \xmark
& --- & --- & --- & ---
& AUC: 0.88--0.95 (valid.), 0.44--0.65 (externos) \\

% VERIFICAR con la ficha de T10: conjuntos de datos y si 0.95 es accuracy en patología (95.13 % en TT1).
Baccouche et al. (2022) \cite{Baccouche2022ResidualNN}
& CBIS-DDSM / INbreast
& \xmark
& Ensemble ResNet
& ---
& \cmark
& \xmark
& \xmark
& --- & --- & --- & ---
& Acc.: 0.95 (patología) \\

\multicolumn{13}{@{}l}{\cellcolor{gray!12}\scriptsize\textit{Aprendizaje automático cuántico fuera de la mamografía}} \\[1pt]

% respaldo: «Model Training accuracy (%) Validation accuracy (%) QCCNN 74.375 70.498 CNN 59.708 58.944 Logistic regression 61.927 60.583 Table 2.» (Xiang2024QuantumCNN, p. 9)
% respaldo: «Model Training accuracy (%) Validation accuracy (%) QCCNN 97.099 96.155 CNN 92.932 92.822 Logistic regression 91.354 91.184 Table 4.» (Xiang2024QuantumCNN, p. 10)
Xiang et al. (2024) \cite{Xiang2024QuantumCNN}
& GBSG / SEER / WDBC (tabulares)
& \cmark\ (sim.)
& QCCNN (PennyLane)
& Capa cuántica conv.
& \xmark
& \xmark
& ---
& Parc. & No & No & No
& Acc. valid.: 0.70--0.96 (clás.: 0.61--0.93) \\

% respaldo: «for ZZFeatureMap, the results were an accuracy of 0.91, precision of 0.87, recall of 0.96, specificity of 0.85, and F1-score of 0.91; and for PauliFeatureMap, the model achieved a consistent performance of 0.96 across all metrics, including accuracy, precision, recall, specificity, and F1-score.» (ElHafidi2025, p. 9)
% respaldo: «The radial basis function (RBF) kernel provided the best results, with an accuracy of 0.95, precision of 0.98, recall of 0.92, specificity of 0.98, and F1-score of 0.95.» (ElHafidi2025, p. 9)
Toufah et al. (2025) \cite{ElHafidi2025}
& Cáncer de pulmón (tabular)
& \cmark\ (sim.)
& QSVM (Qiskit)
& Z / ZZ / Pauli (sin cadenas)
& \xmark
& \xmark
& ---
& Sí & No & ND & No
& Acc.: 0.91--0.96 (RBF: 0.95) \\

% respaldo: «Model: Support Vector Machine (SVM) with a precomputed kernel, implemented using the scikit-learn library with default parameters.» (Parigi2026QImR, p. 6)
% respaldo: «Quantum kernels on average perform binary classification with the same accuracy as the classical linear kernel, but with exponentially lower storage cost.» (Parigi2026QImR, p. 7)
Parigi et al. (2026) \cite{Parigi2026QImR}
& MNIST / F-MNIST (16$\times$16)
& \cmark\ (sim.)
& Kernels cuánticos / SVM
& TNR / FRQI / NEQR / QPIE
& \xmark
& \cmark\ (Gram)
& ---
& Sí & Sí & Sí & No
& Acc.: 0.94--0.99 (lineal: 0.96--0.99) \\

\multicolumn{13}{@{}l}{\cellcolor{gray!12}\scriptsize\textit{Aprendizaje automático cuántico en mamografía}} \\[1pt]

% respaldo: «AUC 0.70 0.58 0.77 0.77 0.74 0.69» (Azevedo2022QuantumTransfer, p. 12)
% respaldo: «Accuracy 0.67 0.84 0.84 0.83 0.83 0.84» (Azevedo2022QuantumTransfer, p. 12)
Azevedo et al. (2022) \cite{Azevedo2022QuantumTransfer}
& BCDR
& \cmark\ (sim. + QPU)
& QTL (PennyLane)
& Angle embedding
& \xmark
& \xmark
& \xmark
& Parc. & No & No & No
& AUC: 0.77 (clás. TL: 0.58); acc.: 0.84 en ambos \\

\midrule
\multicolumn{13}{@{}l}{\cellcolor{blue!8}\scriptsize\textit{Presente trabajo}} \\[1pt]

\rowcolor{blue!8}
\textbf{Presente trabajo}
& \textbf{CBIS-DDSM}
& \cmark\ (\textbf{sim. statevector})
& \textbf{MLP sobre $\langle Z_i\rangle$ (Qiskit)}
& \textbf{C4: ZZ / C5: Pauli (Y, ZZ)}
& \cmark
& \cmark
& \cmark
& Sí & Sí & Sí & No
& % fuente: H-042; Code/results/7_test.csv (análisis main)
  AUC masas: 0.530 / 0.581 (C1: 0.651); calcs.: 0.633 / 0.652 (C1: 0.769) \\

\bottomrule
\end{tabularx}

\vspace{4pt}
\scriptsize
\textbf{Notas:}
\cmark\ = cumple la condición;\quad
\xmark\ = no la cumple;\quad
--- = no aplica;\quad
ND = el trabajo no lo dice;\quad
Parc. = parcialmente.
P1--P4: paridad, incertidumbre, partición de prueba y simulabilidad (sección~\ref{sec:revision_critica}); no aplican a los trabajos sin componente cuántica.
En el presente trabajo, los valores de C4 y C5 van separados por una barra; C1 es la condición sin transformación, con el mismo clasificador.
\end{table}
\end{landscape}
